\documentclass[10pt]{article}

\usepackage[margin=1in]{geometry}
\usepackage{amsmath,amssymb,amsthm}
\usepackage{mathpartir}
\usepackage{booktabs}
\usepackage{array}
\usepackage{microtype}
\usepackage[hidelinks]{hyperref}

% ---------------------------------------------------------------- macros
\newcommand{\ov}[1]{\overline{#1}}
\newcommand{\kw}[1]{\mathbf{#1}}          % keywords
\newcommand{\ksf}[1]{\mathsf{#1}}         % kinds / entries
\newcommand{\This}{\ksf{This}}
\newcommand{\thh}{\hat\theta}
\newcommand{\subst}[1]{[#1]}
\newcommand{\emptytel}{\cdot}
\newcommand{\wf}{\;\ksf{ok}}
\newcommand{\entails}{\Vdash}
\newcommand{\Eval}{\Downarrow}
\newcommand{\dom}{\ksf{dom}}
\newcommand{\clo}[2]{\langle #1,\, #2\rangle}
\newcommand{\unitv}{()}
\newcommand{\smallcode}[1]{\texttt{\small #1}}

\theoremstyle{plain}
\newtheorem{theorem}{Theorem}[section]
\newtheorem{proposition}[theorem]{Proposition}
\theoremstyle{remark}
\newtheorem{remark}[theorem]{Remark}

\title{\textbf{Core Moss: A Calculus of Scope and Context}}
\author{The Moss project --- working note\thanks{This note distills the
\texttt{ml-modules} design iteration of Moss, recorded in
\href{https://github.com/moss-lang/moss/blob/ml-modules/docs/design/semantics.md}{\texttt{docs/design/semantics.md}}
(revision 9), into a core calculus.
Decision tags such as [D1] refer to entries in that log. The definitions
and the entire metatheory of \S\ref{sec:meta} are mechanized and
machine-checked in Rocq in \texttt{CoreMoss.v} alongside this document;
every theorem reports closed under the global context.}}
\date{July 29, 2026}

\begin{document}
\maketitle

\begin{abstract}
Moss separates \emph{scope}---what a name refers to---from
\emph{context}---whether the thing it refers to is available.
Declarations without definitions introduce \emph{abstract symbols}: types,
values, and functions that are in scope but that nobody has yet provided.
The construct $\kw{assume}$ makes a region of code require such symbols,
and $\kw{bind}$ provides them; a requirement is satisfied only by naming
the exact symbol, never by searching for a provider of the right shape.
We present \textsc{Core Moss}, a small calculus distilling this design:
its syntax, static semantics, and a big-step dynamic semantics in which
the \emph{satisfaction maps} computed during elaboration are exactly the
interpreter's environment. The calculus makes the language's central
phase distinction precise---type and function provisions specialize away
at compile time, and values are the only context that exists at run
time---and records the semantics under which a Moss compiler written in
Moss compiles itself to a byte-identical fixpoint. The metatheory,
including type soundness, is machine-checked in Rocq.
\end{abstract}

\section{Introduction}\label{sec:intro}

Most languages resolve two questions with one mechanism: what a name
means, and whether the thing it means is at hand. Moss keeps them apart.
\emph{Scope} is resolved lexically, from imports and enclosing
declarations. \emph{Context} is a separate discipline: a top-level
declaration without a definition---$\kw{type}\;t;$, $\kw{val}\;v:\tau;$,
$\kw{fn}\;f(\ldots):\tau;$---introduces an \emph{abstract symbol}, a need
that is perfectly nameable but not yet provided. A region of code
declares which needs it has with $\kw{assume}$, and code that has the
goods supplies them with $\kw{bind}$.

Two principles govern everything else.

\begin{enumerate}
\item \textbf{Satisfaction is by symbol identity [D1].} A requirement is
met only by an enclosing $\kw{assume}$ of \emph{that} symbol or a
$\kw{bind}$ of \emph{that} symbol. There is no structural matching, no
search of the scope for a provider of the right type, and no specificity
ranking. Method dispatch obeys the same discipline: its lookup key is the
pair (receiver type, method symbol), and it demands exactly one hit
---type-\emph{directed}, never type-\emph{searched}.
\item \textbf{Providing is kind-directed [D2].} Providing a \emph{type}
chooses a representation, statically, and drives monomorphization.
Providing a \emph{function} supplies code together with the context it
captured at the bind site. Providing a \emph{value} supplies runtime
data. After elaboration, types and functions disappear into
specialization; values are the only context that exists at run time.
\end{enumerate}

There is deliberately no more than this. Requirement sets double as the
generics mechanism: a declaration nested in $\kw{assume}$ blocks is
implicitly parameterized by everything they name, and a bracket
application such as $\smallcode{Range[T=NameId]}$ instantiates those
parameters, with anything left unbracketed inherited from the
surrounding $\kw{assume}$s rather than left open [D21, D22]. The design
is validated
empirically: a compiler for Moss, written in Moss, compiles its own
source to a fixpoint---the second and third generations of the compiler
are byte-for-byte identical WebAssembly modules
(\href{https://github.com/moss-lang/moss/blob/ml-modules/docs/implementation/selfhosting.md}{\texttt{docs/implementation/selfhosting.md}}).

This note boils the language down to a core calculus, \textsc{Core
Moss}. Section~\ref{sec:example} shows the surface language on one
example; \S\ref{sec:calculus} gives the core's syntax, statics, and
dynamics; \S\ref{sec:meta} states its metatheory;
\S\ref{sec:full} accounts for everything the boiling removed.

\section{Moss in one example}\label{sec:example}

\begin{figure}[t]
\begin{verbatim}
type Elem;                          # a need: an abstract type symbol

unit Empty;

assume Elem {
  val zero: Elem;                   # a val need; its type mentions Elem
  fn .or_zero(): Elem;              # a detached-method need; receiver chosen
                                    #   where the method is provided
  type Full Elem;                   # a tag, implicitly parameterized by Elem
  type Slot Empty | Full;           # a nominal type over a union

  assume zero {
    fn Slot.coalesce(): Elem {      # attached: defined code, needs {Elem, zero}
      match this { Slot u => match u { Empty => zero, Full e => e } }
    }
  }

  assume zero, Slot.or_zero {
    fn demo(s: Slot): Elem { s.or_zero() }
  }
}

unit O;

fn main() {
  bind Elem = O;                    # a type: static, chooses a representation
  bind zero = O;                    # a val: dynamic, evaluated once, here
  bind Slot.or_zero = Slot.coalesce;# a method: code plus its captured context
  demo(Slot (Full (O)));            # Full injects into the union Empty | Full
}
\end{verbatim}
\caption{A complete Moss program exercising the core: needs of all three
kinds, telescoped \texttt{assume}s, implicit parameterization, nominal
data with injection, and the three flavors of \texttt{bind}.}
\label{fig:example}
\end{figure}

Figure~\ref{fig:example} is a complete program. \smallcode{Elem} is an
abstract type; \smallcode{zero} is an abstract val whose declared type
mentions it, so its declaration must sit inside
$\kw{assume}\;\smallcode{Elem}$---requirements are \emph{telescoped}:
an item may be assumed only where the symbols its own signature mentions
are already available [D20]. The tag \smallcode{Full} is declared under
the same assumption, so \smallcode{Full} is implicitly parameterized by
\smallcode{Elem}; that is the whole of generics [D21]. Type identity is
structural over (declaration, static bindings), so once
$\kw{bind}\;\smallcode{Elem=O}$ is in force, \smallcode{Full} means the
same instantiated type everywhere downstream [D18].

\smallcode{.or\_zero} is a \emph{detached method}: a need whose receiver
type is supplied where the method is provided, so one symbol can be
provided at many receivers, each provision keyed by the pair (receiver,
method) [D36]. \smallcode{demo} requires that key
(\smallcode{Slot.or\_zero}); the call \smallcode{s.or\_zero()} resolves by
forward-inferring the receiver's type and demanding exactly one provision
at that key. The provider bound in \smallcode{main} is the \emph{attached}
method \smallcode{Slot.coalesce}---ordinary defined code that can see its
receiver---and the bind captures the provider's own requirements
(\smallcode{zero}, a runtime value) at the bind site. The three binds in
\smallcode{main} illustrate the kind-directed split [D2]: the type bind is
purely static, the val bind evaluates a value, and the method bind pairs
static code with captured dynamic context.

\section{The core calculus}\label{sec:calculus}

\textsc{Core Moss} strips Moss to the constructs that carry the design.
Modules, named contexts, functors, attached methods, records, tuples,
booleans, loops, and mutation are all omitted here and reconstructed in
\S\ref{sec:full}; none of them adds semantic content to the core.

\subsection{Syntax}

\begin{figure}[t]
\[
\begin{array}{llcl}
\text{nominal refs} & \nu & ::= & N[\theta] \\
\text{types} & \tau & ::= & \unitv \mid t \mid \This \mid
  \nu_1 \vee \cdots \vee \nu_k \quad (k \ge 0) \\
\text{applications} & \theta & ::= & [\,\ov{t = \tau}\,] \\[4pt]
\text{items} & \iota & ::= & t \mid v \mid f \mid \tau.m[\theta] \\
\text{telescopes} & \Delta & ::= & \emptytel \mid \Delta, \iota \\[4pt]
\text{declarations} & d & ::= &
    \kw{type}\;t
  \mid \kw{type}\;N\;\tau
  \mid \kw{val}\;v : \tau
  \mid \kw{fn}\;f(\ov{x : \tau}) : \tau \\
& & \mid & \kw{fn}\;f(\ov{x : \tau}) : \tau \;\{\,e\,\}
  \mid \kw{fn}\;.m(\ov{x : \tau}) : \tau
  \mid \kw{assume}\;\Delta\;\{\,\ov{d}\,\} \\[4pt]
\text{expressions} & e & ::= &
    x \mid \unitv \mid v
  \mid N[\theta]\;e
  \mid f[\theta](\ov{e})
  \mid e.m(\ov{e}) \\
& & \mid & \kw{match}\;e\;\{\,\ov{N\;x \Rightarrow e}\,\}
  \mid \kw{let}\;x = e;\; e
  \mid \kw{bind}\;\beta;\; e \\[2pt]
\text{binds} & \beta & ::= &
    t = \tau \mid v = e \mid f = f' \mid \tau.m[\theta] = f' \\[4pt]
\text{programs} & P & ::= & \ov{d}
  \qquad\text{with } \kw{fn}\;\ksf{main}() : \unitv\;\{\,e\,\} \text{ at the top level} \\[4pt]
\text{values} & w & ::= & \unitv \mid N\,w \\
\text{contexts in force} & \Phi & ::= & \langle A ;\, \sigma \rangle
  \qquad A \subseteq_{\mathrm{fin}} \{\iota\},\;
  \sigma : t \rightharpoonup \tau \\
\text{runtime contexts} & \delta & ::= &
  \ov{\iota \mapsto w} \cup \ov{\iota \mapsto \clo{f}{\delta}}
\end{array}
\]
\caption{Syntax of \textsc{Core Moss}. Metavariables: $t$ ranges over
abstract type symbols, $N$ over nominal tag names, $v$ over val symbols,
$f$ over fn symbols, $m$ over method symbols, $x$ over variables. We
write $N$ for $N[\,]$, $\bot$ for the empty union ($k=0$; Moss spells it
\texttt{|}), and $e_1; e_2$ for $\kw{let}\;x = e_1;\,e_2$ with $x$ fresh.
$\kw{unit}\;N$ abbreviates $\kw{type}\;N\;\unitv$, with the value $N$
short for $N\,\unitv$ and the pattern $N$ for $N\,x$ with $x$ unused.}
\label{fig:syntax}
\end{figure}

Figure~\ref{fig:syntax} gives the syntax. Symbols of the four kinds are
drawn from disjoint countable sets; there are no literal expressions and,
importantly, \emph{no first-class functions}: an $f$ position names a
symbol, so a function value never flows as data---code travels only
through $\kw{bind}$, paired with its captured context. Union members are
syntactically nominal references with distinct heads [D15], and a tag's
payload may itself be a union (that is how the example's \smallcode{Slot}
and the library's \smallcode{Bool} are built).

A \emph{telescope} $\Delta$ is an ordered requirement list; the
declarations under $\kw{assume}\;\Delta$ have $\Delta$ appended to their
requirements. A telescope may mention one \emph{key} twice---the key of a
type, val, or fn item is the symbol itself, and of a method item the pair
(receiver type, method symbol)---because flattening named contexts
concatenates their members. Formation \emph{merges} such duplicates by
unifying the colliding bindings [D43]; the merging judgment of
Figure~\ref{fig:statics-aux} makes this precise. Elaborating a program yields a global signature $\Sigma$
assigning every symbol its kind, its declared signature, and the
concatenation of its enclosing telescopes:
\[
\begin{array}{c}
\Sigma(t) = \ksf{type} \qquad
\Sigma(N) = \ksf{tag}[\Delta]\,\tau \qquad
\Sigma(v) = \ksf{val}[\Delta]\,\tau \\[2pt]
\Sigma(f) = \ksf{fn}[\Delta](\ov{x{:}\tau} \to \tau)\;\{e\}^? \qquad
\Sigma(m) = \ksf{mth}[\Delta](\ov{x{:}\tau} \to \tau)
\end{array}
\]
A method signature may mention the distinguished abstract type $\This$,
the receiver's type; in the core a method's arguments are understood as
$(\This, \ov{\tau})$, i.e.\ every provider receives the receiver as its
first argument (\S\ref{sec:full} explains how surface Moss spells this).

\paragraph{Contexts in force.}
At each program point the statics tracks $\Phi = \langle A; \sigma
\rangle$: the set $A$ of items available (assumed or bound) and the
static substitution $\sigma$ accumulated by type binds. Assumed abstract
types are in $A$ but not in $\dom(\sigma)$; bound ones are in both.
$\sigma$ acts on types and items homomorphically, and we write $\tau
\equiv \tau'$ for equality of $\sigma$-images (unions up to reordering);
this is the applicative identity rule---the same declaration under the
same static bindings is the same type everywhere [D18]. Item
availability, written $\iota \in \Phi$, means $\sigma\iota' = \sigma\iota$
for some $\iota' \in A$; comparing after $\sigma$ is Moss's key
canonicalization [D43], and $\sigma$ carries the entries produced by
telescope merging (Figure~\ref{fig:statics-aux}) alongside those of type
binds, so
two spellings of one merged atom are one item everywhere the statics
compares.

\paragraph{Nothing dangling.}
A reference must leave none of its target's requirements unresolved, but
it need not answer them all in brackets: what the enclosing
$\kw{assume}$s already supply, they supply. Given $\theta$ with its
domain $\dom(\theta)$ (the type symbols it binds) contained in the type
items of $\Delta$, its \emph{completion} $\thh$ extends $\theta$ with the
identity on the remaining type items of $\Delta$---the pass-through
reading, under which a requirement not mentioned in brackets is
re-exported as the very same symbol. The judgment
$\Sigma;\Phi \vdash \theta \entails \Delta$
(Figure~\ref{fig:statics-aux}) checks that this leaves nothing dangling:
every type requirement is either instantiated by $\theta$ or already in
force, and every other requirement is available after substitution. All
four reference forms---type references, tag construction, defined-fn
calls, and method provisions---reuse this one judgment.

The two extremes are the common ones: a fully applied
$\smallcode{Range[T=NameId]}$, and a bare $\smallcode{Range}$ that
substitutes nothing and thus refers to the very $\smallcode{T}$ its
surroundings assume (the \texttt{inner.moss} shared-$\smallcode{T}$
idiom). The mixed case is admitted on the same footing: inside
$\kw{assume}\;\smallcode{U}$, a reference $f[\smallcode{T=Char}]$ to an
$f$ requiring $\{\smallcode{T},\smallcode{U}\}$ lets $\smallcode{U}$ pass
through. What [D22] rules out is not this but a partially instantiated
reference \emph{travelling}: after elaboration every requirement of every
reference is pinned, either by $\theta$ or to an in-force symbol, so
nothing downstream ever has to adapt a half-applied provider to a
differently-shaped need.

\subsection{Static semantics}

\begin{figure}[t]
\centering
\textbf{Keys and unification}
\[
\begin{array}{r@{\;}c@{\;}l@{\qquad}r@{\;}c@{\;}l}
\ksf{key}(t) &=& t &
\ksf{key}(f) &=& f \\
\ksf{key}(v) &=& v &
\ksf{key}(\tau.m[\theta]) &=& (\tau, m) \\[3pt]
\multicolumn{6}{l}{\ksf{keys}(\Delta) \;=\; \{\,\ksf{key}(\iota) \mid \iota \in \Delta\,\}}\\[3pt]
\multicolumn{6}{l}{\ksf{unify}(\iota_1, \iota_2) \;=\;
  \text{the most general }\sigma_u\text{ with }
  \sigma_u\iota_1 = \sigma_u\iota_2,\ \text{if one exists}}
\end{array}
\]

\textbf{Merging [D43]}\quad
\fbox{$\Delta \Downarrow \langle \hat\Delta;\, \sigma_\Delta \rangle$}
\begin{mathpar}
\inferrule[M-Nil]{ }{\emptytel \Downarrow \langle \emptytel;\, \emptyset \rangle}

\inferrule[M-Fresh]{\Delta \Downarrow \langle \hat\Delta;\, \sigma_\Delta \rangle \\
  \ksf{key}(\sigma_\Delta\iota) \notin \ksf{keys}(\sigma_\Delta\hat\Delta)}
  {\Delta, \iota \Downarrow \langle \hat\Delta, \iota;\, \sigma_\Delta \rangle}

\inferrule[M-Merge]{\Delta \Downarrow \langle \hat\Delta;\, \sigma_\Delta \rangle \\
  \iota' \in \hat\Delta \\
  \ksf{key}(\sigma_\Delta\iota') = \ksf{key}(\sigma_\Delta\iota) \\
  \ksf{unify}(\sigma_\Delta\iota',\, \sigma_\Delta\iota) = \sigma_u}
  {\Delta, \iota \Downarrow \langle \hat\Delta;\, \sigma_u \circ \sigma_\Delta \rangle}
\end{mathpar}

\textbf{Satisfaction}\quad \fbox{$\Sigma;\Phi \vdash \theta \entails \Delta$}
\begin{mathpar}
\inferrule[S-Tel]{\Delta \Downarrow \langle \hat\Delta;\, \sigma_\Delta \rangle \\
  \Sigma;\Phi \vdash \theta \entails_1 \hat\Delta \\
  \sigma(\thh\,t) = \sigma(\thh\,\tau) \;\;(\forall\, t \mapsto \tau \in \sigma_\Delta)}
  {\Sigma;\Phi \vdash \theta \entails \Delta}

\inferrule[S-Nil]{ }{\Sigma;\Phi \vdash \theta \entails_1 \emptytel}

\inferrule[S-Inst]{\Sigma;\Phi \vdash \theta \entails_1 \hat\Delta \\
  t \in \dom(\theta) \\ \Sigma;\Phi \vdash \theta(t) \wf}
  {\Sigma;\Phi \vdash \theta \entails_1 \hat\Delta, t}

\inferrule[S-Pass]{\Sigma;\Phi \vdash \theta \entails_1 \hat\Delta \\
  t \notin \dom(\theta) \\ t \in \Phi}
  {\Sigma;\Phi \vdash \theta \entails_1 \hat\Delta, t}

\inferrule[S-Item]{\Sigma;\Phi \vdash \theta \entails_1 \hat\Delta \\
  \iota \text{ not a type symbol} \\ \thh\,\iota \in \Phi \\
  \sigma(\thh\,t) = \sigma(t) \;\;(\forall t \in \ksf{deps}(\iota))}
  {\Sigma;\Phi \vdash \theta \entails_1 \hat\Delta, \iota}
\end{mathpar}

\textbf{Type formation}\quad \fbox{$\Sigma;\Phi \vdash \tau \wf$}
\begin{mathpar}
\inferrule[K-Unit]{ }{\Sigma;\Phi \vdash \unitv \wf}

\inferrule[K-Need]{t \in \Phi}{\Sigma;\Phi \vdash t \wf}

\inferrule[K-Tag]{\Sigma(N) = \ksf{tag}[\Delta]\,\tau_0 \\
  \Sigma;\Phi \vdash \theta \entails \Delta}
  {\Sigma;\Phi \vdash N[\theta] \wf}

\inferrule[K-Union]{\Sigma;\Phi \vdash \nu_i \wf \quad (\forall i) \\
  \sigma\nu_1, \ldots, \sigma\nu_k \text{ have distinct heads}}
  {\Sigma;\Phi \vdash \nu_1 \vee \cdots \vee \nu_k \wf}
\end{mathpar}

\textbf{Subtyping (injection only)}\quad \fbox{$\Sigma;\Phi \vdash \tau \le \tau'$}
\begin{mathpar}
\inferrule[Sub-Eq]{\tau \equiv \tau'}{\Sigma;\Phi \vdash \tau \le \tau'}

\inferrule[Sub-Inj]{\sigma\nu = \sigma\nu_j \text{ for some } j}
  {\Sigma;\Phi \vdash \nu \le \nu_1 \vee \cdots \vee \nu_k}
\end{mathpar}
\caption{Merging, satisfaction, type formation, and subtyping.
\emph{Merging} [D43]: a method item's key is the pair (receiver
\emph{type}, method symbol), dropping its application $\theta$; keys are
compared after $\sigma_\Delta$, and $\sigma$ acts on a telescope
itemwise. Keying a method on the receiver's \emph{type} is the same [D18]
identity that dispatch is keyed by [D62], and for the same reason: keying on
its head would give $\smallcode{Slot[Elem=Int].get}$ and
$\smallcode{Slot[Elem=Char].get}$ one key, so a telescope needing the
method at two instantiations would try to identify $\smallcode{Int}$ with
$\smallcode{Char}$ and be rejected. D43's own example still merges: there
the receiver is shared and it is the method's own application that
differs. $\ksf{unify}$ is ordinary first-order syntactic
unification---symmetric, neither occurrence primary---so the two ways
formation can fail are the two ways an mgu can fail to exist: a
\emph{clash} of distinct concrete heads (say
$\smallcode{A.gimme[Foo=Eof]}$ with $\smallcode{A.gimme[Foo=Comma]}$), and
a \emph{cyclic} equation $t \doteq \tau$ with $t$ free in $\tau$ (the
occurs check; D43 is silent on cyclic merges, a finding of the
mechanization). $\sigma_u \circ \sigma_\Delta$ is the idempotent
extension: the new equations are applied to the old ranges.
\emph{Satisfaction} is judged against the merged telescope
(\textsc{S-Tel}): totality is per equivalence class---$\dom(\theta)$ covers
the type items of $\hat\Delta$---and $\theta$ must respect the merge
equations, so instantiating one spelling of a merged atom instantiates them
all. In \textsc{S-Inst}/\textsc{S-Pass}/\textsc{S-Item}, $\dom(\theta)$
must be contained in the type items of $\hat\Delta$; $\sigma$ is the
substitution component of $\Phi$; $\ksf{deps}(\iota)$ is the set of type
symbols in $\iota$'s own telescope. \textsc{S-Item}'s last premise is
\emph{coherence}: a val or fn item names no types, so without it nothing
stops $f[T{=}\ksf{Bool}]$ from being satisfied by an ambient $v : T$ that
was provided under $\kw{bind}\;T{=}\ksf{Int}$. (Found while mechanizing
preservation; method items carry their types and need no such premise.)}
\label{fig:statics-aux}
\end{figure}

\begin{figure}[t]
\centering
\fbox{$\Sigma;\Phi;\Gamma \vdash e : \tau$}
\begin{mathpar}
\inferrule[T-Var]{\Gamma(x) = \tau}{\Sigma;\Phi;\Gamma \vdash x : \tau}

\inferrule[T-Unit]{ }{\Sigma;\Phi;\Gamma \vdash \unitv : \unitv}

\inferrule[T-Sub]{\Sigma;\Phi;\Gamma \vdash e : \tau' \\
  \Sigma;\Phi \vdash \tau' \le \tau}
  {\Sigma;\Phi;\Gamma \vdash e : \tau}

\inferrule[T-Val]{\Sigma(v) = \ksf{val}[\Delta]\,\tau \\ v \in \Phi}
  {\Sigma;\Phi;\Gamma \vdash v : \sigma\tau}

\inferrule[T-Tag]{\Sigma(N) = \ksf{tag}[\Delta]\,\tau_0 \\
  \Sigma;\Phi \vdash \theta \entails \Delta \\
  \Sigma;\Phi;\Gamma \vdash e : \sigma\thh\tau_0}
  {\Sigma;\Phi;\Gamma \vdash N[\theta]\;e : \sigma(N[\thh])}

\inferrule[T-Need]{\Sigma(f) = \ksf{fn}[\Delta](\ov{x{:}\tau} \to \tau_r) \\
  f \in \Phi \\
  \Sigma;\Phi;\Gamma \vdash e_i : \sigma\tau_i \quad (\forall i)}
  {\Sigma;\Phi;\Gamma \vdash f(\ov{e}) : \sigma\tau_r}

\inferrule[T-Call]{\Sigma(f) = \ksf{fn}[\Delta](\ov{x{:}\tau} \to \tau_r)\,\{e_b\} \\
  \Sigma;\Phi \vdash \theta \entails \Delta \\
  \Sigma;\Phi;\Gamma \vdash e_i : \sigma\thh\tau_i \quad (\forall i)}
  {\Sigma;\Phi;\Gamma \vdash f[\theta](\ov{e}) : \sigma\thh\tau_r}

\inferrule[T-Meth]{\Sigma(m) = \ksf{mth}[\Delta](\ov{x{:}\tau} \to \tau_r) \\
  \Sigma;\Phi;\Gamma \vdash e_0 : \tau_0 \\
  \tau'.m[\theta] \in \Phi \\ \sigma\tau' = \sigma\tau_0 \\
  \rho = \sigma \circ \subst{\This \mapsto \sigma\tau_0} \circ \thh \\
  \Sigma;\Phi;\Gamma \vdash e_i : \rho\tau_i \quad (\forall i)}
  {\Sigma;\Phi;\Gamma \vdash e_0.m(\ov{e}) : \rho\tau_r}

\inferrule[T-Match]{\Sigma;\Phi;\Gamma \vdash e_0 : \tau_0 \\
  \sigma\tau_0 = N_1[\theta_1] \vee \cdots \vee N_k[\theta_k] \\
  \Sigma(N_i) = \ksf{tag}[\Delta_i]\,\tau_i^{p} \\
  \Sigma;\Phi;\Gamma, x_i : \theta_i\tau_i^{p} \vdash e_i : \tau \quad (\forall i)}
  {\Sigma;\Phi;\Gamma \vdash
    \kw{match}\;e_0\;\{\,\ov{N\;x \Rightarrow e}\,\} : \tau}

\inferrule[T-Let]{\Sigma;\Phi;\Gamma \vdash e_1 : \tau_1 \\
  \Sigma;\Phi;\Gamma, x : \tau_1 \vdash e_2 : \tau}
  {\Sigma;\Phi;\Gamma \vdash \kw{let}\;x = e_1;\;e_2 : \tau}
\end{mathpar}
\caption{Expression typing. \textsc{T-Meth} carries the side condition
that the provision at the key $(\sigma\tau_0, m)$ --- the receiver's
\emph{type}, exactly as \S\ref{sec:intro} says, not merely its head ---
is unique in $\Phi$ \emph{modulo $\sigma$}: provisions are keyed by the
receiver's (declaration, static bindings) identity [D18] compared through
the merged substitution, so two spellings of one merged atom are one
provision rather than an ambiguity [D62]. Keying on the head alone is
adequate under Moss's unique-nominal-per-instantiation idiom [D51] but
unsound in the raw calculus, a repair owed to the mechanization. In
\textsc{T-Match}, a nominal reference is identified with its singleton
union, arms are in bijection with the heads $N_1, \ldots, N_k$, and $k =
0$ types $\kw{match}\;e_0\;\{\}$ at any $\tau$.}
\label{fig:typing}
\end{figure}

\begin{figure}[t]
\centering
\fbox{$\Sigma;\Phi;\Gamma \vdash e : \tau$ (continued: $\kw{bind}$)}
\begin{mathpar}
\inferrule[T-BindTy]{\Sigma(t) = \ksf{type} \\ t \neq \This \\ t \notin \Phi \\
  \Sigma;\Phi \vdash \tau' \wf \\
  \Sigma;\langle A \cup \{t\};\, \sigma[t \mapsto \sigma\tau'] \rangle;\Gamma \vdash e : \tau \\
  \Sigma;\Phi \vdash \tau \wf}
  {\Sigma;\langle A; \sigma \rangle;\Gamma \vdash \kw{bind}\;t = \tau';\;e : \tau}

\inferrule[T-BindVal]{\Sigma(v) = \ksf{val}[\Delta_v]\,\tau_v \\
  \Sigma;\Phi \vdash \emptyset \entails \Delta_v \\
  \Sigma;\Phi;\Gamma \vdash e_1 : \sigma\tau_v \\
  \Sigma;\langle A \cup \{v\}; \sigma \rangle;\Gamma \vdash e : \tau}
  {\Sigma;\langle A; \sigma \rangle;\Gamma \vdash \kw{bind}\;v = e_1;\;e : \tau}

\inferrule[T-BindFn]{\Sigma(f) = \ksf{fn}[\Delta_f]\,S \\
  \Sigma(f') = \ksf{fn}[\Delta']\,S'\,\{e_b\} \\
  \Sigma;\Phi \vdash \emptyset \entails \Delta_f \\
  \Sigma;\Phi \vdash \emptyset \entails \Delta' \\
  \sigma S' = \sigma S \\
  \Sigma;\langle A \cup \{f\}; \sigma \rangle;\Gamma \vdash e : \tau}
  {\Sigma;\langle A; \sigma \rangle;\Gamma \vdash \kw{bind}\;f = f';\;e : \tau}

\inferrule[T-BindMeth]{\Sigma(m) = \ksf{mth}[\Delta_m](\ov{x{:}\tau} \to \tau_r) \\
  \Sigma;\Phi \vdash \tau_\kappa \wf \\
  \Sigma;\Phi \vdash \theta \entails \Delta_m \\
  \rho = \sigma \circ \subst{\This \mapsto \sigma\tau_\kappa} \circ \thh \\
  \Sigma(f') = \ksf{fn}[\Delta']\,
    ((y_0{:}\tau_0', \ov{y{:}\tau'}) \to \tau')\,\{e_b\} \\
  \Sigma;\Phi \vdash \emptyset \entails \Delta' \\
  \sigma\tau_0' = \sigma\tau_\kappa \\
  \sigma\ov{\tau'} = \rho\ov{\tau} \\ \sigma\tau' = \rho\tau_r \\
  \Sigma;\langle A \cup \{\tau_\kappa.m[\theta]\}; \sigma \rangle;\Gamma \vdash e : \tau}
  {\Sigma;\langle A; \sigma \rangle;\Gamma \vdash
    \kw{bind}\;\tau_\kappa.m[\theta] = f';\;e : \tau}
\end{mathpar}
\caption{Typing the three kinds of provision (method binds shown
separately from plain fn binds). $S$ ranges over signatures
$\ov{x{:}\tau} \to \tau$, compared pointwise; in \textsc{T-BindFn} the
provided $f$ must be abstract (no body in $\Sigma$). Signature equality
is checked \emph{after} the substitutions in force, and it is an error,
not a unification opportunity, if abstract symbols remain unmatched
[D27].}
\label{fig:binds}
\end{figure}

\begin{figure}[t]
\centering
\textbf{Telescopes}\quad \fbox{$\Sigma;\Delta_0 \vdash \Delta \wf$}
\quad (write $\Phi_{\Delta} = \langle \widehat{\Delta_0, \Delta};\,
\sigma_{\Delta_0, \Delta} \rangle$, the merged context)
\begin{mathpar}
\inferrule[W-Nil]{ }{\Sigma;\Delta_0 \vdash \emptytel \wf}

\inferrule[W-Type]{\Sigma;\Delta_0 \vdash \Delta \wf \\ \Sigma(t) = \ksf{type} \\
  t \neq \This}
  {\Sigma;\Delta_0 \vdash \Delta, t \wf}

\inferrule[W-Sym]{\Sigma;\Delta_0 \vdash \Delta \wf \\
  \Sigma(s) = \ksf{val}[\Delta_s]\,\tau \text{ or } \ksf{fn}[\Delta_s]\,S \\
  \Sigma;\Phi_{\Delta} \vdash \emptyset \entails \Delta_s}
  {\Sigma;\Delta_0 \vdash \Delta, s \wf}

\inferrule[W-Meth]{\Sigma;\Delta_0 \vdash \Delta \wf \\
  \Sigma(m) = \ksf{mth}[\Delta_m]\,S \\
  \Sigma;\Phi_{\Delta} \vdash \tau_\kappa \wf \\
  \Sigma;\Phi_{\Delta} \vdash \theta \entails \Delta_m}
  {\Sigma;\Delta_0 \vdash \Delta, \tau_\kappa.m[\theta] \wf}
\end{mathpar}

\textbf{Declarations}\quad \fbox{$\Sigma;\Delta_0 \vdash d \wf$}
\quad (write $\Phi_0 = \langle \Delta_0; \emptyset \rangle$)
\begin{mathpar}
\inferrule[D-Need]{ }{\Sigma;\Delta_0 \vdash \kw{type}\;t \wf}

\inferrule[D-Tag]{\Sigma;\Phi_0 \vdash \tau \wf}
  {\Sigma;\Delta_0 \vdash \kw{type}\;N\;\tau \wf}

\inferrule[D-Val]{\Sigma;\Phi_0 \vdash \tau \wf}
  {\Sigma;\Delta_0 \vdash \kw{val}\;v : \tau \wf}

\inferrule[D-Fn]{\Sigma;\Phi_0 \vdash \tau_i \wf \quad (\forall i) \\
  \Sigma;\Phi_0 \vdash \tau_r \wf}
  {\Sigma;\Delta_0 \vdash \kw{fn}\;f(\ov{x{:}\tau}) : \tau_r \wf}

\inferrule[D-FnDef]{\Sigma;\Phi_0 \vdash \tau_i \wf \quad (\forall i) \\
  \Sigma;\Phi_0 \vdash \tau_r \wf \\
  \Sigma;\Phi_0;\ov{x : \tau} \vdash e : \tau_r}
  {\Sigma;\Delta_0 \vdash \kw{fn}\;f(\ov{x{:}\tau}) : \tau_r\;\{e\} \wf}

\inferrule[D-Meth]{\Sigma;\langle \Delta_0 \cup \{\This\}; \emptyset \rangle
    \vdash \tau_i \wf \quad (\forall i) \\
  \Sigma;\langle \Delta_0 \cup \{\This\}; \emptyset \rangle \vdash \tau_r \wf}
  {\Sigma;\Delta_0 \vdash \kw{fn}\;.m(\ov{x{:}\tau}) : \tau_r \wf}

\inferrule[D-Assume]{\Sigma;\Delta_0 \vdash \Delta \wf \\
  \Sigma;\Delta_0, \Delta \vdash d_i \wf \quad (\forall i)}
  {\Sigma;\Delta_0 \vdash \kw{assume}\;\Delta\;\{\ov{d}\} \wf}
\end{mathpar}
\caption{Telescope and declaration formation. A program $P$ is
well-formed, $\vdash P \wf$, when its declarations induce a
function $\Sigma$ (each symbol declared once, with the concatenation of
its enclosing telescopes recorded), $\Sigma;\emptytel \vdash d \wf$ for
each top-level $d$, and $\Sigma(\ksf{main}) =
\ksf{fn}[\emptytel](\to \unitv)\{e\}$. \textsc{W-Sym} is the telescoping
rule [D20]: an item may be assumed only after the items its own
declaration needs. Formation additionally requires the merge of
Figure~\ref{fig:statics-aux} to succeed---$\Delta_0, \Delta \Downarrow
\langle \widehat{\Delta_0, \Delta}; \sigma_{\Delta_0, \Delta} \rangle$---and
every context the rules consult is that merged context, so a declaration's
body is checked with the merge equations in force.}
\label{fig:decls}
\end{figure}

Figures~\ref{fig:statics-aux}--\ref{fig:decls} give the statics.
Everything specific to Moss is visible in a handful of premises:

\begin{itemize}
\item \emph{Availability, not search.} \textsc{T-Val}, \textsc{T-Need},
\textsc{T-Meth}, and \textsc{S-Item} test $\iota \in \Phi$---finite-set
membership modulo $\sigma$. No rule quantifies over $\Sigma$ looking for
a provider whose type fits [D1].
\item \emph{Two ways to call.} An \emph{abstract} fn is called because it
is available (\textsc{T-Need}); its own telescope was checked where it
entered $\Phi$. A \emph{defined} fn is called by satisfying its telescope
on the spot (\textsc{T-Call}); the premise $\theta \entails \Delta$ is
where implicit parameterization is discharged, and $\sigma\thh$ is the
monomorphizing substitution.
\item \emph{Dispatch is directed.} \textsc{T-Meth} forward-infers the
receiver's type and demands the unique provision at $(\sigma\tau_0, m)$
--- the receiver's type itself, compared after the substitutions in
force. The provision's own bracket bindings $\theta$
are what interpret the method's signature at the call site---ergonomics
ride along with the provision instead of being discovered [D36].
\item \emph{Duplicate keys merge, consistently.} Telescope formation
unifies colliding bindings (\textsc{M-Merge}) instead of rejecting them,
erring only when two distinct concrete heads would be identified [D43].
The equations land in $\sigma$, so downstream nothing changes: availability
and dispatch already compare modulo $\sigma$, and \textsc{S-Tel} makes a
use site instantiate merged symbols consistently---``binding one binds
them all''.
\item \emph{Injection is the only subtyping.} \textsc{Sub-Inj} lets a
member into a union that lists it; there is no width or depth subtyping
and no union-to-union coercion [D17].
\item \emph{Binds never infer.} In Figure~\ref{fig:binds} every signature
comparison happens after $\sigma$; a bind whose signatures still mention
distinct abstract symbols simply fails, which forces the surface
program's characteristic ``bind the types first, then the functions''
order [D27]. A method bind may carry $\theta$ on its \emph{left-hand}
side---that is how a method declared once (say \smallcode{.get} with
element type \smallcode{Elem}) is provided at one receiver at a time,
with the substitution local to that provision [D61].
\end{itemize}

\subsection{Dynamic semantics}

The dynamics is defined on \emph{elaborated} programs. The typing
derivation resolves, at each site, exactly which in-force item satisfied
each requirement; elaboration records these choices:

\begin{itemize}
\item each defined-fn call $f[\theta]^{\varsigma}(\ov{e})$ and each bind
of a defined fn carries a \emph{satisfaction map} $\varsigma$ sending
every non-type item of the callee's (resp.\ provider's) telescope to the
item of $A$ that satisfied it;
\item each method call $e_0.m^{\iota_0}(\ov{e})$ carries the provision
$\iota_0 \in A$ that \textsc{T-Meth} selected.
\end{itemize}

\noindent
This is precisely the compilation model: after elaboration nothing ever
looks at an $\kw{assume}$ or resolves a name again, and the interpreter's
environment \emph{is} the explicit context structure elaboration
produced. A runtime context $\delta$ maps val items to values and
fn/method items to closures $\clo{f'}{\delta'}$---a (code,
captured-context) pair, the dynamic residue of [D2]. Given $\varsigma$,
we write $\delta \circ \varsigma$ for the context $\{\iota \mapsto
\delta(\varsigma\,\iota)\}$, defined on the non-type items of the
relevant telescope: the callee's requirements, re-keyed under the
callee's own names. Type items have no runtime content and are dropped.
Figure~\ref{fig:dynamics} gives the rules.

\begin{figure}[t]
\centering
\fbox{$\delta;\gamma \vdash e \Eval w$}
\begin{mathpar}
\inferrule[E-Unit]{ }{\delta;\gamma \vdash \unitv \Eval \unitv}

\inferrule[E-Var]{\gamma(x) = w}{\delta;\gamma \vdash x \Eval w}

\inferrule[E-Val]{\delta(v) = w}{\delta;\gamma \vdash v \Eval w}

\inferrule[E-Tag]{\delta;\gamma \vdash e \Eval w}
  {\delta;\gamma \vdash N[\theta]\,e \Eval N\,w}

\inferrule[E-Let]{\delta;\gamma \vdash e_1 \Eval w_1 \\
  \delta;\gamma[x \mapsto w_1] \vdash e_2 \Eval w}
  {\delta;\gamma \vdash \kw{let}\;x = e_1;\;e_2 \Eval w}

\inferrule[E-Match]{\delta;\gamma \vdash e_0 \Eval N_j\,w' \\
  \delta;\gamma[x_j \mapsto w'] \vdash e_j \Eval w}
  {\delta;\gamma \vdash \kw{match}\;e_0\;\{\,\ov{N\;x \Rightarrow e}\,\} \Eval w}

\inferrule[E-Call]{\Sigma(f) = \ksf{fn}[\Delta](\ov{x{:}\tau} \to \tau_r)\,\{e_b\} \\
  \delta;\gamma \vdash e_i \Eval w_i \quad (\forall i) \\
  (\delta \circ \varsigma);\, [\ov{x \mapsto w}] \vdash e_b \Eval w}
  {\delta;\gamma \vdash f[\theta]^{\varsigma}(\ov{e}) \Eval w}

\inferrule[E-Need]{\delta(f) = \clo{f'}{\delta'} \\
  \Sigma(f') = \ksf{fn}[\Delta'](\ov{y{:}\tau} \to \tau_r)\,\{e_b\} \\
  \delta;\gamma \vdash e_i \Eval w_i \quad (\forall i) \\
  \delta';\, [\ov{y \mapsto w}] \vdash e_b \Eval w}
  {\delta;\gamma \vdash f(\ov{e}) \Eval w}

\inferrule[E-Meth]{\delta(\iota_0) = \clo{f'}{\delta'} \\
  \Sigma(f') = \ksf{fn}[\Delta']((y_0{:}\tau_0', \ov{y{:}\tau'}) \to \tau')\,\{e_b\} \\
  \delta;\gamma \vdash e_0 \Eval w_0 \\
  \delta;\gamma \vdash e_i \Eval w_i \quad (\forall i) \\
  \delta';\, [y_0 \mapsto w_0,\, \ov{y \mapsto w}] \vdash e_b \Eval w}
  {\delta;\gamma \vdash e_0.m^{\iota_0}(\ov{e}) \Eval w}

\inferrule[E-BindTy]{\delta;\gamma \vdash e \Eval w}
  {\delta;\gamma \vdash \kw{bind}\;t = \tau';\;e \Eval w}

\inferrule[E-BindVal]{\delta;\gamma \vdash e_1 \Eval w_1 \\
  \delta[v \mapsto w_1];\gamma \vdash e \Eval w}
  {\delta;\gamma \vdash \kw{bind}\;v = e_1;\;e \Eval w}

\inferrule[E-BindFn]{\delta[f \mapsto \clo{f'}{\delta \circ \varsigma}];\gamma
    \vdash e \Eval w}
  {\delta;\gamma \vdash \kw{bind}\;f =^{\varsigma} f';\;e \Eval w}

\inferrule[E-BindMeth]{\delta[\tau_\kappa.m[\theta] \mapsto
    \clo{f'}{\delta \circ \varsigma}];\gamma \vdash e \Eval w}
  {\delta;\gamma \vdash \kw{bind}\;\tau_\kappa.m[\theta] =^{\varsigma} f';\;e \Eval w}
\end{mathpar}
\caption{Big-step dynamics of elaborated \textsc{Core Moss}. A program
runs as $\emptyset;\emptyset \vdash e_{\ksf{main}} \Eval \unitv$. Note
what is absent: \textsc{E-BindTy} is erasure, applications $\theta$ never
reach a value, and no rule consults a type.}
\label{fig:dynamics}
\end{figure}

Three points deserve emphasis. First, \textsc{E-Call} runs a defined
body under $\delta \circ \varsigma$---the caller's provisions,
restricted and renamed to the callee's requirement list. This is
context \emph{propagation}: nothing dynamic is looked up by name; the
plumbing was fixed by elaboration. Second, \textsc{E-BindFn} and
\textsc{E-BindMeth} capture $\delta \circ \varsigma$ at the bind
site---a provided function is a pair of static code and dynamic context,
so the same abstract symbol bound in two places may carry different
captured data. Third, \textsc{E-BindVal} evaluates its right-hand side
exactly once, at the bind [D26]; val provisions are call-by-value data,
not thunks.

\section{Metatheory}\label{sec:meta}

\begin{proposition}[Resolution is search-free]\label{prop:nosearch}
Every context premise in Figures~\ref{fig:statics-aux}--\ref{fig:decls}
is a membership test $\iota \in \Phi$ in a finite set modulo the
idempotent substitution $\sigma$; no rule quantifies over the program
signature $\Sigma$. Consequently type checking is syntax-directed (up to
the placement of \textsc{T-Sub}) and elaboration---the assignment of
satisfaction maps $\varsigma$ and provisions $\iota_0$---is a partial
function of the program: at every use site there is at most one way to
satisfy each requirement---up to the substitution in force, under which
the spellings of a merged atom are interchangeable [D43, D62]---and typing
fails rather than choosing.
\end{proposition}

\begin{theorem}[Soundness]\label{thm:sound}
Say $\delta \models \Phi$ when $\delta$ is defined on every non-type item
of $\Phi$ and each entry agrees with the item's declared signature under
the substitutions in force (values at declared types, closures at matched
signatures with $\delta' \models$ their own telescopes), and $\gamma
\models \Gamma$ pointwise. If\, $\vdash P \wf$,\;
$\Sigma;\Phi;\Gamma \vdash e : \tau$,\; $\delta \models \Phi$,\;
$\gamma \models \Gamma$, and $\delta;\gamma \vdash e \Eval w$, then
$w$ has type $\sigma\tau$; and in the instrumented semantics that signals
$\ksf{err}$ whenever a lookup $\gamma(x)$, $\delta(\iota)$, or a match
arm is missing, no evaluation of $e$ signals $\ksf{err}$
(\emph{context safety}: a well-typed program never reaches for a
provision that is not there). In particular a terminating run of a
well-formed program yields $\unitv$.
\end{theorem}

The proof is a routine induction on the evaluation derivation with
inversion on typing, with a substitution lemma for $\sigma\thh$ and a
weakening lemma for $\delta \circ \varsigma$. The metatheory is fully
machine-checked in Rocq (\texttt{CoreMoss.v}, checked in CI): this
theorem and every proposition here report \emph{closed under the global
context}, via a fuel-indexed safety invariant that carries a ground
instantiation of the abstract type symbols per closure.  The
mechanization additionally carries an executable fueled interpreter:
the program of Figure~\ref{fig:example}, encoded in the core, computes
to $\unitv$ by \texttt{vm\_compute}.  Designing and proving the
invariant is what surfaced the repairs below. (One further
mechanization-level deviation: the mechanized merge is \emph{oriented}---the
earlier occurrence of a key is kept as the representative and the unifier
emits entries mapping later symbols onto it---and the unifier's output is
verified at formation rather than trusted, so no property of the unifier
enters the metatheory; the symmetric presentation above is recovered up to
choice of representative.) The theorem's empirical
counterpart is blunter: the self-hosted compiler elaborates and runs
its own 19-module source under exactly these rules, to a byte-identical
fixpoint.

\begin{proposition}[Phase separation]\label{prop:phase}
Evaluation never inspects a type: the relation $\Eval$ is invariant
under erasing every $\theta$, every $\kw{bind}\;t = \tau$, and the type
items of every telescope. Moreover, in every $\delta$ reachable from a
well-formed program, the code component of every closure is determined
by the elaboration annotations alone---only the captured val data varies
at run time. Hence a whole-program specializer may resolve every
\textsc{E-Need} and \textsc{E-Meth} to a direct call and inline every
type binding, leaving val provisions as the only context crossing
runtime edges [D2].
\end{proposition}

\paragraph{What mechanizing the proofs caught.}
The repairs below, most surfaced while designing the preservation
invariant and one by the merging session, are already folded into the
figures above. (i)~Dispatch must key on the
receiver's \emph{type}: an earlier draft of \textsc{T-Meth} keyed
provisions by the receiver's head, which is adequate under Moss's
unique-nominal-per-instantiation idiom [D51] but lets a provision bound
at $\smallcode{Slot[Int]}$ serve a $\smallcode{Slot[Char]}$ receiver,
the $\This$-substitution lying about the result.
(ii)~Satisfaction needs \emph{coherence} (\textsc{S-Item}'s last
premise): a val or fn requirement names no types, so nothing else
prevents $f[T{=}\ksf{Bool}]$ from being satisfied by an ambient
$v : T$ that was provided under $\kw{bind}\;T{=}\ksf{Int}$.
(iii)~Type binds must be fresh (\textsc{T-BindTy}'s $t \notin \Phi$):
under substitution-based statics, shadowing an in-force type symbol
conflates its two generations; surface Moss's lexical shadowing [D25]
is instead the province of environment-passing semantics, and real
binds provide symbols not yet in force, so the restriction does not
bite the corpus.
(iv)~The bind's result type may not escape its scope
(\textsc{T-BindTy}'s final premise): without it, subsumption under the
extended $\sigma$ retypes the body at $t$ itself, and the bound symbol
leaks into a type the outside cannot read --- refuted mechanically by a
proof agent before it was a premise.
(v)~$\rho$ must compose with the $\This$-substitution \emph{outermost}
(\textsc{T-Meth}/\textsc{T-BindMeth}): an earlier draft applied
$[\This \mapsto \sigma\tau_0]$ innermost, where $\thh$ can rewrite
inside the receiver --- refuted through a method item that mentions a
symbol of its own application's domain; the corrected order subsumes
two stopgap coherence premises this document briefly carried.
(vi)~Subsumption must target well-formed types and \textsc{T-Match}
must know its scrutinee's $\sigma$-image is well-formed: regularity the
declarative rules otherwise lose, and which the soundness invariant
leans on.
(vii)~Cyclic merges must fail (the occurs check in
$\ksf{unify}$): [D43] specifies merging as deterministic unification but is
silent on the cyclic case, and without the check a context could require
$t \equiv \smallcode{Full[Elem=}t\smallcode{]}$, a requirement no
instantiation grounds.
(viii)~$\This$ is \emph{reserved}: nothing else stopped a program from
binding the receiver symbol itself and pre-capturing every receiver
substitution in scope --- the only repair that refuted
Theorem~\ref{thm:sound} outright (a well-formed program whose
$\ksf{main}$ evaluates to $\ksf{err}$, checked by \texttt{vm\_compute}).
In surface Moss this is unutterable, since \texttt{This} is a keyword
[D3]; the core forgot to encode the lexer's opinion as a side
condition.  Two of the findings above were subsequently tested against
the bootstrap compiler [D62]: (i) is a \emph{live} hole --- the
bootstrap keys provisions by the receiver's head symbol and accepts a
program that types a $\smallcode{Char}$ payload as
$\smallcode{Int}$ --- and (ii) is latent, unreachable only because v0
does not yet implement bracket applications at call sites. The repair
[D62] is now part of the calculus: provisions are keyed by the receiver's
(declaration, static bindings) identity, and---what merging makes
practical---compared modulo the merged $\sigma$, so re-keying does not
fracture provisions across the spellings of one atom.

Proposition~\ref{prop:phase} is the hypothesis behind the compilation
pipeline: the monomorphizing WebAssembly backend is partial evaluation
of Figure~\ref{fig:dynamics} with respect to its static parts, and
contains no resolution logic of its own. One caveat survives from the
full language: because type binds drive specialization, a recursive
function that rebinds a type on its recursive path demands infinitely
many specializations; Moss accepts such programs at elaboration time and
bounds the instantiation graph with a depth backstop, the sole check
permitted after elaboration [D30, D16].

\section{From the core to Moss}\label{sec:full}

Everything the surface language adds is scoping convenience, sugar over
the core, or an orthogonal extension. Table~\ref{tab:full} is the
accounting.

\begin{table}[h]
\centering\small
\begin{tabular}{@{}p{0.30\linewidth}p{0.64\linewidth}@{}}
\toprule
\textbf{Moss} & \textbf{In the core} \\
\midrule
Modules, \texttt{import}/\texttt{use}/\texttt{as} &
Pure name management with no context effect; a module is a namespace,
not a structure. $\alpha$-resolution folds the module graph into the
core's flat symbol space [D8--D10, D44, D56]. \\[3pt]
\texttt{context} declarations &
Names for telescopes; assuming one assumes its members, flattened
recursively [D19]. Duplicate keys \emph{merge} by unification, erring only
when two distinct concrete types would be identified [D43]---and the core
now carries exactly this rule at telescope formation
(Figure~\ref{fig:statics-aux}), so flattening is pure concatenation. \\[3pt]
Functors &
A named, signature-checked block of binds, applied with
$\kw{bind}\;F;$ --- application inlines the block, so it elaborates to a
run of core binds [D55]. \\[3pt]
Attached methods, \texttt{this}/\texttt{This} &
$\kw{fn}\;T.m$ is a defined fn found by name in the receiver's own
module; a detached provision is found by key [D61]. Only an attached
method can see its receiver, so providers wrap nominal types [D54]; the
core flattens this by giving every method provider the receiver as its
first argument---the convention both compilers already use internally. \\[3pt]
Records and tuples &
Structural products, orthogonal to the context mechanism; nominality
comes only from tags wrapping them [D15]. \\[3pt]
\texttt{Bool}, \texttt{if}, \texttt{while}/\texttt{loop}, \texttt{var} &
\texttt{Bool} is a library tag over two units and \texttt{if} is its
match [D45]. Loops are primitive in Moss because tail-call elimination
is deliberately \emph{not} semantics [D49]; the core, having no loops,
uses recursion. Local \texttt{var} is a reassignable slot; shared
mutable state goes through the \texttt{Cell} library [D31]. \\[3pt]
Literals &
None, in Moss or the core: values enter programs through the
environment (\texttt{char::a}, \texttt{zero}, \texttt{first\_arg()})
[D4, D5]. \\[3pt]
Operators &
Reserved; they will desugar to detached methods provided per receiver
type (\texttt{a + b} $\equiv$ \texttt{a.add(b)}) [D33, D47]. \\[3pt]
Entry and the primitive context &
The primitive context is \texttt{Wasi}; \texttt{Std} is a Moss library
over it, installed by a functor, and \texttt{main} may assume any subset
of it [D37, D52]. The core's \texttt{main} assumes nothing and binds its
own provisions. \\
\bottomrule
\end{tabular}
\caption{What was boiled away, and where it went.}
\label{tab:full}
\end{table}

\section{Related work}\label{sec:related}

Core Moss sits between two traditions and refuses the defining move of
each. From type classes~\cite{wadler-blott} and their descendants---the
implicit calculus~\cite{implicit-calculus}, Scala implicits, OCaml's
modular implicits~\cite{modular-implicits}---it takes implicit
program-wide plumbing of capabilities, but rejects resolution by
type-directed \emph{search}: a Moss requirement names one symbol, and
coherence is by construction rather than by global instance discipline.
From ML modules it takes the vocabulary---a \texttt{context} is a
signature, the provisions in force are a structure, functors map
structures to structures---together with applicative type
identity~\cite{leroy-applicative}, but it dissolves the structure as a
runtime value: provisions are delivered along lexically-checked edges,
and only val data survives to run time. The nearest ancestor of
$\kw{bind}$ for vals is implicit parameters~\cite{implicit-params},
dynamic scoping with static types; Core Moss extends the idea uniformly
to types and code, where ``dynamic'' scoping becomes static
specialization. The consistent-merging rule for contexts [D43]---in the core, the
\textsc{M-Merge} rule of telescope formation---is a deterministic
congruence over requirement sets;
its nearest published relative is the record-like treatment of
intersection types of Xu, Huang, and Oliveira~\cite{type-difference}.

\begin{thebibliography}{9}\small

\bibitem{wadler-blott}
P.~Wadler and S.~Blott.
\newblock How to make \emph{ad-hoc} polymorphism less \emph{ad hoc}.
\newblock In \emph{POPL}, 1989.
\newblock \url{https://doi.org/10.1145/75277.75283}

\bibitem{implicit-params}
J.~R. Lewis, J.~Launchbury, E.~Meijer, and M.~B. Shields.
\newblock Implicit parameters: Dynamic scoping with static types.
\newblock In \emph{POPL}, 2000.
\newblock \url{https://doi.org/10.1145/325694.325708}

\bibitem{implicit-calculus}
B.~C. d.~S. Oliveira, T.~Schrijvers, W.~Choi, W.~Lee, and K.~Yi.
\newblock The implicit calculus: A new foundation for generic programming.
\newblock In \emph{PLDI}, 2012.
\newblock \url{https://doi.org/10.1145/2254064.2254070}

\bibitem{modular-implicits}
L.~White, F.~Bour, and J.~Yallop.
\newblock Modular implicits.
\newblock In \emph{ML Family Workshop}, 2014.
\newblock \url{https://doi.org/10.4204/EPTCS.198.2}

\bibitem{leroy-applicative}
X.~Leroy.
\newblock Applicative functors and fully transparent higher-order modules.
\newblock In \emph{POPL}, 1995.
\newblock \url{https://doi.org/10.1145/199448.199476}

\bibitem{type-difference}
H.~Xu, X.~Huang, and B.~C. d.~S. Oliveira.
\newblock Making a type difference: Subtraction on intersection types as
  generalized record operations.
\newblock \emph{Proc.\ ACM Program.\ Lang.}, 7(POPL), 2023.
\newblock \url{https://doi.org/10.1145/3571224}

\end{thebibliography}

\end{document}
